\chapter{Giới thiệu}
\label{ch:gioithieu}

\section{Đặt vấn đề}

Khi thao tác trên điện thoại, người dùng nhiều lúc dừng lại giữa chừng vì không rõ bước
tiếp theo phải làm gì. Điều họ cần trong tình huống đó là một câu chỉ dẫn cho đúng bước
trước mắt, chẳng hạn ``Chạm vào biểu tượng chia sẻ ở góc trên bên phải màn hình''. Với dạng
câu hướng dẫn ngắn và chỉ nói về màn hình người dùng đang mở như vậy, đối tượng tiếp nhận
là con người chứ không phải phần mềm.

Dòng nghiên cứu về tác tử giao diện (\emph{GUI agent}) những năm gần đây giải một bài toán
rất gần với bài toán trên, nhưng khác hẳn ở bản chất của đầu ra. Cho cùng một ảnh màn hình và cùng một mục
tiêu, các hệ thống như SeeClick~\cite{seeclick}, OS-Atlas~\cite{osatlas} hay
Aguvis~\cite{aguvis} sinh ra một lệnh máy đọc được, dạng \code{click(872,141)}, để phần
mềm tự bấm thay người. Câu chữ, nếu có xuất hiện, chỉ là biểu diễn trung gian dẫn tới lệnh đó và không được đem chấm.

Luận văn nhận cùng đầu vào đó, nhưng lấy câu chữ làm sản phẩm cuối cùng. Đầu ra là một
câu hướng dẫn cho người đọc, và đó cũng là thứ duy nhất được đem đánh giá. Việc đổi đầu ra
như vậy kéo theo hai vấn đề chưa có lời giải sẵn.

\textbf{Vấn đề thứ nhất là đánh giá.} Một toạ độ dự đoán có thể đem so trực tiếp với toạ
độ mà người dùng đã chạm. Một câu thì không so được với một câu chuẩn duy nhất, bởi cùng
một nút thường có nhiều cách gọi tên đều đúng. Chẳng hạn ``Click on the search bar'' và
``Tap the magnifying glass at the top'' cùng chỉ một phần tử, và không bộ dữ liệu nào liệt kê hết
các cách diễn đạt hợp lệ.

Các thước đo quen dùng đều thất bại ở đúng điểm này, và luận văn đo mức thất bại của từng
thước. So khớp chuỗi ký tự loại nhầm $97{,}5\%$ số câu diễn đạt khác. So vector ngữ nghĩa cho
diện tích dưới đường ROC (AUC) chỉ $0{,}336$ trên $51$ trường hợp viết tay, tức còn kém hơn
đoán ngẫu nhiên. Lý do là vector ngữ nghĩa đo độ gần về \emph{chủ đề}, trong khi thứ cần
đo lại là hai câu có nói về \emph{cùng một phần tử} hay không. Ở phép thử này, cặp
\code{inbox} và \code{outbox} được chấm $0{,}76$, cao hơn hẳn cặp \code{search} và \code{magnifying
glass} chỉ được $0{,}55$, dù cặp sau mới là cặp cùng chỉ một nút. Chi tiết ở
Mục~\ref{sec:thuoc-gay}.

\textbf{Vấn đề thứ hai nằm ở dạng lỗi mà mô hình sinh câu hay mắc.} Đọc tay $40$ trường
hợp sinh sai cho thấy việc mô hình bịa ra một nút không có trên màn hình chỉ chiếm khoảng
$0$ đến $2\%$. Phần lớn lỗi là câu \emph{mơ hồ}. Mô hình gọi ``biểu tượng tìm kiếm'' trong khi màn hình
có tới hai biểu tượng tìm kiếm. Một câu như vậy không sai về mặt sự thật, nhưng không dùng được với người đọc vì câu ấy không phân biệt được phần tử cần chạm với các phần tử giống nó bên cạnh. Quan sát này
đặt lại trọng tâm của cải tiến. Phần đáng sửa là \emph{tính phân biệt} của câu, còn việc chống bịa đặt gần như không còn khoảng trống để cải thiện.

Hai vấn đề này quyết định bố cục của luận văn. Một phần dành cho việc dựng và kiểm chứng
một thước đo dùng được; phần còn lại dành cho việc thiết kế và huấn luyện một mô hình sinh
câu ít mơ hồ hơn.

\section{Động lực nghiên cứu}

\subsection{Động lực khoa học}

Câu hỏi ``một mô tả có đủ để bên kia nhận ra đúng đối tượng hay không'' không mới. Dòng
sinh biểu thức quy chiếu (\emph{referring expression generation}) trong thị giác máy
tính đã đặt đúng câu hỏi đó từ năm $2016$. Mao và cộng sự~\cite{mao2016} chấm một mô tả
bằng cách xem một mô hình hiểu (\emph{listener}) có chỉ đúng đối tượng không; Luo và
Shakhnarovich~\cite{luo2017} cũng như Yu và cộng sự~\cite{yu2017} còn đưa mô hình hiểu
đó vào thẳng vòng huấn luyện. Nguyên tắc này được kế thừa nguyên vẹn ở đây, và đề tài
không nhận là mình nghĩ ra nó.

Khoảng trống nằm ở miền giao diện. Widget Captioning~\cite{widgetcap}, công trình gần
nhất trong miền này, nêu rõ vấn đề hai phần tử trông giống hệt nhau nhưng vẫn huấn
luyện bằng hàm mất mát entropy chéo thuần, tức là không có thành phần nào ép mô tả phải
phân biệt. Như vậy, tính phân biệt đã được nhận diện như một vấn đề trong miền
giao diện nhưng chưa được đưa vào \emph{mục tiêu huấn luyện}.

Ở phía đánh giá, việc kiểm chứng một phát biểu bằng một nguồn có cấu trúc \emph{bên
ngoài} thay vì bằng chính niềm tin của mô hình sinh là nguyên tắc đứng sau
ALOHa~\cite{aloha}. Hướng đi ở đây bám theo nguyên tắc đó, song có một khác biệt làm bảo đảm ấy yếu đi. ALOHa đối chiếu với một danh sách tĩnh, còn ở đây đối chiếu với một
\emph{mô hình đã học}, nên độ tin cậy của thước bị chặn trên bởi độ tin cậy của mô hình
đó. Vì vậy một phần đáng kể của luận văn dành cho việc đo đạc chính dụng cụ đo
trước khi dùng nó để kết luận bất cứ điều gì.

Cuối cùng, có một cảnh báo nhắm vào chính thước đề xuất.
Jandial và cộng sự~\cite{groundersunderstand} cho thấy độ chính xác mà các bộ chuẩn hiện
nay báo cáo, đo trên \emph{một} câu tốt nhất cho mỗi phần tử, là con số thổi phồng so
với năng lực thật. Một tác tử đối kháng chuyên chế câu phá mô hình đạt tỉ lệ thành công tới
$84\%$ trên giao diện máy tính để bàn. Con số đó là \emph{năng suất của một phép tìm kiếm
đối kháng}, không phải tỉ lệ trượt khi người dùng diễn đạt khác đi, còn trường hợp giao diện di động thì họ để ngỏ. Mục~\ref{sec:phepA} đo trực tiếp phần cảnh báo áp
được vào đây, Mục~\ref{sec:hanche} nêu phần còn hở.

\subsection{Động lực ứng dụng}

Nhu cầu thực tế cho tác vụ này đã có. GuideMe~\cite{guideme} là một hệ thống hỗ trợ người cao
tuổi tự học dùng điện thoại thông minh. Người dùng hỏi, hệ thống nhìn màn hình và trả về
hướng dẫn từng bước. Hệ thống đó gọi một mô hình thương mại cỡ lớn qua mạng và không
huấn luyện gì; chất lượng được đánh giá qua một nghiên cứu người dùng gồm $18$ người,
không có thước định lượng nào đi kèm.

Từ đó có hai khoảng trống mang tính ứng dụng. Thứ nhất, chưa có thước định lượng
nào cho phép so sánh hai hệ thống sinh hướng dẫn giao diện với nhau một cách tái lập
được. Thứ hai, việc phải gửi ảnh màn hình lên dịch vụ ngoài là một rào cản thật. Ảnh
màn hình điện thoại chứa tin nhắn, số dư tài khoản và danh bạ. Một mô hình cỡ ba tỉ tham
số chạy được tại máy người dùng giải quyết đúng rào cản đó, và đó là lý do luận văn chọn
kích cỡ mô hình như vậy thay vì chọn mô hình lớn nhất có thể.

\section{Mục đích nghiên cứu}

Đề tài đặt ra bốn mục tiêu.

\begin{enumerate}[leftmargin=1.6em]
\item Xây dựng một thước đo cho bài toán sinh hướng dẫn giao diện không dựa vào việc so
sánh với một câu chuẩn duy nhất, và \emph{kiểm chứng} thước đó bằng các phép đo độc lập
với ngưỡng đóng băng trước khi chạy.

\item Chuyển bộ dữ liệu AndroidControl~\cite{androidcontrol} thành tập giám sát cho tác
vụ này một cách hoàn toàn tự động, không thuê người dán nhãn và không dùng nhãn do mô
hình khác sinh ra.

\item Thiết kế và huấn luyện một mô hình thị giác và ngôn ngữ cỡ nhỏ cho tác vụ, với một thành phần nhắm vào dạng lỗi đã đo được là câu mơ hồ. Thành phần này bắt mô hình
mô tả phần tử cần chạm theo một dạng có cấu trúc trước khi viết câu.

\item Đăng ký trước thiết kế so sánh cùng luật đọc kết quả cho cả bốn kết cục có thể
xảy ra, để kết luận không phụ thuộc vào việc nhìn thấy số trước hay sau.
\end{enumerate}

\section{Một số nghiên cứu liên quan}

\subsection{Định vị phần tử giao diện và tác tử giao diện}

AndroidControl~\cite{androidcontrol} cung cấp $15.283$ chuỗi thao tác do người thật thực
hiện trên thiết bị Android, mỗi bước kèm ảnh màn hình, toạ độ chạm và một câu hướng dẫn
do chính người thao tác viết. AndroidInTheWild~\cite{aitw} là nguồn của quy ước dung sai
$14\%$ cạnh màn được dùng rộng rãi để tính một dự đoán toạ độ là trúng hay trượt.

Các mô hình định vị phần tử như SeeClick~\cite{seeclick}, OS-Atlas~\cite{osatlas} và
UGround~\cite{uground} ánh xạ một biểu thức văn bản sang một điểm trên màn hình. Luận
văn dùng một mô hình thuộc nhóm này làm \textbf{dụng cụ đo}, không phải làm đường cơ sở
để so sánh, và vì vậy phải báo cáo sai số đo được cùng thành phần dữ liệu huấn luyện của mô hình ấy trước khi phụ thuộc vào nó (Mục~\ref{sec:congA} và Mục~\ref{sec:hanche}).

Việc đưa một danh sách phần tử có cấu trúc vào đầu vào cũng đã phổ biến.
Chính nhóm tác giả AndroidControl tinh chỉnh mô hình với danh sách phần tử suy từ cây
trợ năng làm đầu vào và nói rõ tác tử của họ không dùng trực tiếp ảnh màn
hình~\cite{androidcontrol}; Mind2Web~\cite{mind2web} huấn luyện trên tập ứng viên có
nhiễu; Widget Captioning~\cite{widgetcap} và Screen2Words~\cite{screen2words} đặt cây
phần tử ở đầu vào. Đề tài \emph{không} nhận đóng góp nào cho tầng này; nó chỉ xuất
hiện như một nhánh đối chứng.

\subsection{Sinh mô tả có tính phân biệt}

Bài toán sinh một mô tả đủ để nhận ra đúng một đối tượng giữa các đối tượng gây nhiễu là
bài toán sinh biểu thức quy chiếu~\cite{mao2016}. Các công trình sau đó đặt một mô hình
hiểu vào trong vòng huấn luyện để ép tính phân biệt~\cite{luo2017,yu2017}. Trong miền
giao diện, công trình có giám sát gần nhất là Widget Captioning~\cite{widgetcap}, và như
đã nêu, hàm mất mát ở đó là entropy chéo thuần.

Vì dòng nghiên cứu này đã tồn tại từ $2016$, luận văn \textbf{không dùng các từ ``đầu
tiên'', ``mới'' hay ``cơ chế mới''} cho thành phần đề xuất. Phần đóng góp được phát biểu
hẹp lại. Luận văn đưa tính phân biệt vào đích huấn luyện cho bài toán sinh hướng dẫn
giao diện, rồi kiểm xem nó có chuyển giao được sang miền này hay không.

\subsection{Sinh bước trung gian trước khi sinh đáp án}

Việc bắt mô hình sinh một biểu diễn trung gian trước khi sinh đáp án cuối đã được thử
nhiều lần, và các kết quả \emph{không} cùng chiều. Bảng~\ref{tab:tienle} tổng hợp năm
điểm dữ liệu đã tra tận nguồn.

\begin{table}[t]
\caption{Ảnh hưởng của bước trung gian trong các công trình trước. Số dương nghĩa là có
bước trung gian thì tốt hơn.}
\label{tab:tienle}
\centering\small
\renewcommand{\arraystretch}{1.2}
\begin{tabular}{@{}p{4.2cm}p{5.2cm}r@{}}
\toprule
Công trình & Phép so & Chênh lệch \\
\midrule
Aguvis~\cite{aguvis} & bỏ tầng trung gian, đo trên chính AndroidControl mức thấp & $-11{,}4$ \\
Shikra (văn xuôi)~\cite{shikra} & bước trung gian viết bằng văn tự do & $-7{,}4$ \\
Shikra (có toạ độ)~\cite{shikra} & cùng mô hình, bước trung gian có toạ độ & $+5{,}9$ \\
GCoT~\cite{gcot} & huấn luyện sinh định vị trước (LLaVA-7B / 13B) & $+4{,}5$ / $+5{,}8$ \\
CogCoM~\cite{cogcom} & bỏ $70$K dữ liệu chuỗi trung gian khỏi huấn luyện & $+6{,}6$ / $+0{,}2$ / $+0{,}9$ \\
\bottomrule
\end{tabular}
\end{table}

Hai dòng Shikra là cặp có giá trị thông tin cao nhất trong bảng, vì chúng dùng cùng một
mô hình và cùng một bài kiểm, chỉ khác nhau ở \emph{dạng} của bước trung gian. Viết bằng
văn xuôi thì kết quả âm, còn viết có kèm toạ độ thì kết quả dương. Đó là căn cứ để dòng khai báo trong luận văn bắt buộc có cấu trúc và có ô
toạ độ, thay vì để mô hình tự do mô tả bằng lời.

Con số âm hay được nhắc tới của GCoT~\cite{gcot} dễ bị trích một nửa, vì nó nằm ở chế độ ra lệnh cho mô hình \emph{chưa} huấn luyện; khi huấn luyện
theo thứ tự định vị trước thì chính công trình đó cho kết quả dương như trong bảng.
Ngoài ra, hướng của các công trình trên là sinh trung gian để cuối cùng ra \emph{toạ
độ}; ở luận văn này là để ra \emph{câu chữ cho người đọc}.

\subsection{Đánh giá không cần đáp án tham chiếu}

Nguyên tắc đối chiếu văn bản sinh ra với một nguồn có cấu trúc bên ngoài, thay vì với
chính mô hình sinh, là nền của ALOHa~\cite{aloha}. Việc kiểm chứng một thước đo bằng
danh sách các phép bơm lỗi với ngưỡng công bố trước là phương pháp đã được thiết lập
bởi Sai và cộng sự~\cite{sai2021}, và luận văn áp dụng đúng quy trình đó
(Mục~\ref{sec:bomloi}). Có hai cảnh báo liên quan trực tiếp tới hướng đi này. Các đại
lượng tự động thay cho đánh giá của người rất dễ bị tin quá mức~\cite{clark2021}, và các
mô hình định vị giao diện đã được báo cáo là mong manh trước việc đổi cách diễn
đạt~\cite{groundersunderstand}.

\section{Nội dung và phạm vi thực hiện đề tài}

\subsection{Nội dung nghiên cứu}

\begin{itemize}[leftmargin=1.4em]
\item Khảo sát các mô hình thị giác và ngôn ngữ, cùng các kỹ thuật tinh chỉnh hiệu quả
tham số phù hợp với ràng buộc một card đồ hoạ đơn.
\item Khảo sát các thước đo hiện có cho văn bản sinh không có đáp án chuẩn, đo mức độ
thất bại của chúng trên dữ liệu thật, từ đó thiết kế thước thay thế.
\item Xây dựng đường ống tự động chuyển AndroidControl thành tập giám sát, gồm bốn khâu
ghép dữ liệu đa nguồn, chạy OCR, dựng nhãn mô tả từ cây trợ năng và kiểm rò rỉ.
\item Huấn luyện và đánh giá mô hình theo thiết kế so sánh đã đăng ký trước.
\item Phân tích lỗi trên từng bước, khai báo giới hạn kèm số đo loại trừ hoặc xác nhận từng giới hạn.
\end{itemize}

\subsection{Phạm vi nghiên cứu}

\textbf{Về bài toán}, luận văn chỉ xét trường hợp \emph{một màn hình}. Đầu vào là một ảnh
chụp màn hình kèm một mục tiêu, còn đầu ra là hướng dẫn cho bước kế tiếp. Trường hợp nhiều
màn hình liên tiếp nằm ngoài phạm vi.

\textbf{Về thước đo}, luận văn chỉ chấm một khía cạnh hẹp và có tên rõ ràng. Câu hỏi được chấm là câu sinh ra
có xác định đúng phần tử cần chạm hay không. Độ trôi chảy, độ dài, mức hữu ích với người đọc
nằm ngoài thước.

\textbf{Về dữ liệu}, thực nghiệm chạy trên AndroidControl, giao diện tiếng Anh. Luận văn
không dựng bộ dữ liệu tiếng Việt; phần tiếng Việt nếu có chỉ ở mức minh hoạ định tính.

\textbf{Về mô hình}, tất cả các nhánh đều tinh chỉnh cùng một mô hình gốc
Qwen2.5-VL-3B-Instruct~\cite{qwen25vl} bằng QLoRA~\cite{qlora} trên một card A100 đơn.
Việc huấn luyện mô hình từ đầu, cũng như so sánh nhiều mô hình gốc khác nhau, nằm ngoài
phạm vi.

\section{Đóng góp của luận văn}

Luận văn có \textbf{ba đóng góp}. Đóng góp thứ nhất thuộc về đo lường và đã hoàn tất.
Đóng góp thứ hai là một kết quả thực nghiệm đã đăng ký trước nhưng \emph{không hoàn tất
được} vì ngân sách máy. Đóng góp thứ ba là một chẩn đoán ở mức từng bước, chỉ ra giới hạn
thật sự của thành phần được đề xuất. Ba phần độc lập nhau về lập luận, nên
phần chưa hoàn tất không kéo theo hai phần còn lại.

\paragraph{Đóng góp thứ nhất, về thước đo.} Luận văn đề xuất và kiểm chứng
\textbf{executability}, thước đo cho khả năng xác định phần tử của một câu hướng dẫn
giao diện. Phần kiểm chứng gồm sáu khối, mỗi khối là một phép đo chứ không phải một lập
luận. Khối thứ nhất là một cổng chặn đặt trước về sai số của mô hình định vị, và cổng này
đạt với sai số trung vị $0{,}7\%$ bề ngang màn so với ngưỡng $3\%$. Khối thứ hai chọn luật
xác định trúng dựa trên \emph{sàn} đo được, vì luật dung sai quy ước vẫn cho điểm $84{,}3\%$
số câu chỉ sai nút trong khi luật ô Voronoi chỉ cho $2{,}8\%$. Khối thứ ba đo hai đầu của
dải điểm. Trần của thước được đo bằng cách đưa chính câu do người viết đi chấm, cho
$75{,}7\%$ với khoảng tin cậy $[74{,}1; 77{,}3]$ trên đủ $4.463$ bước, còn sàn được đo bằng
ba nhánh đối chứng, trong đó mô hình định vị vẫn trúng $12{,}0\%$ số bước khi mọi câu bị
thay bằng một câu chung chung không nói phần tử nào, và chỉ còn trúng $6{,}1\%$ khi thay
bằng câu thật nhưng mô tả một màn hình khác. Khối thứ tư là một bộ mười phép bơm lỗi với
ngưỡng đóng băng trước khi chạy. Khối thứ năm là phép viết lại câu, cho thấy thước bền
trước việc đổi động từ và đảo trật tự câu với hiệu ròng $+0{,}35$ điểm trên $1.139$ câu viết
lại, mà vẫn phạt đúng khi mệnh đề vị trí bị lấy đi. Khối thứ sáu là một phép kiểm chéo
bằng mô hình định vị thứ hai không dùng AndroidControl khi huấn luyện, và phép kiểm này
giữ được $94\%$ độ lớn của chênh lệch dùng làm phép so đối chiếu.

\paragraph{Đóng góp thứ hai, về đường ống dữ liệu và một kết quả đã đăng ký trước.} Luận văn
xây dựng hoàn toàn tự động một tập giám sát $64.567$ bước, kiểm phép ghép dữ liệu đa nguồn
bằng một đối chứng lệch chủ ý ($48\%$ so với $20\%$), và xác nhận không có tác vụ nào trùng
giữa tập huấn luyện và tập kiểm. Trên nền đó, thành phần đề xuất ``mô tả phân biệt trước, phát
ngôn sau'' được huấn luyện và chấm cùng nhánh đối chứng của nó. Kết quả thu được là nhánh
xử lý thấp hơn nhánh nền $1{,}93$ điểm ở một hạt giống, con số nằm trong dải \emph{không kết luận
được} đã khoá từ trước, và hạt giống thứ hai bị huỷ nên kết cục đóng lại ở đó. Chặng huấn
luyện tiếp theo bằng mục tiêu ưu tiên ở tầng khai báo đạt $60{,}0\%$, cao nhất trong mọi
điểm lưu của đề tài, nhưng phần thuộc riêng mục tiêu ấy chỉ là $0{,}63$ điểm trên nhánh
đối chứng, tức \textbf{$78\%$ mức tăng thuộc về đối chứng}; cả hai đều là kết quả một hạt
giống.

\paragraph{Đóng góp thứ ba, về chẩn đoán ở mức từng bước.} Cắt quần thể theo việc ô khai báo
do chính mô hình sinh ra có trỏ đúng phần tử hay không cho thấy cơ chế của thành phần đề
xuất \emph{vẫn hoạt động đúng như giả thuyết}, nhưng bị chặn bởi độ chính xác của chính ô
khai báo đó. Trên $60{,}6\%$ số bước mà mô hình khai báo đúng, nó hơn nhánh chỉ sinh câu
$8{,}83$ điểm và vượt cả trần câu người viết của nhóm đó; trên phần còn lại nó thấp hơn $11{,}12$ điểm. Hai chiều gần như triệt tiêu nhau, nên cộng lại trên toàn quần thể chỉ còn
một chênh lệch $+0{,}94$ điểm không có ý nghĩa thống kê. Đóng góp này còn gồm hai phần phụ.
Thứ nhất là một hệ số chuyển đổi đo được từ ``đúng khai báo'' sang ``đúng executability''
bằng $0{,}43$, cho phép quy một mức cải thiện ở tầng khai báo thành dự báo về điểm đầu ra và
kiểm lại dự báo đó bằng phép đo sau. Thứ hai là lý do đo được vì sao một biến thể đã đăng ký
không dựng nổi dữ liệu huấn luyện. Khi gọi sai tên, mô hình hoặc gọi đúng phần tử bằng một
tên khác, hoặc nhìn sang một vùng khác hẳn của màn hình, chứ không lẫn giữa hai nút cạnh
nhau như biến thể đó giả định.

\paragraph{Về mức độ hoàn thành.} Đại lượng chính của thiết kế đòi trung bình hai hạt giống
cho mỗi nhánh; nhánh xử lý chỉ có một. Luận văn vì vậy \textbf{không} phát biểu rằng thành
phần đề xuất có tác dụng hay gây hại, và cũng không thay đại lượng đó bằng một phép so khác
đã có số. Mọi kết quả của nhánh xử lý mang nhãn \emph{thăm dò} ở từng chỗ chúng xuất hiện,
và Bảng~\ref{tab:trangthai} ghi trạng thái cuối cùng của cả bảy nhánh đã đăng ký kèm lý do
cho từng nhánh không chạy.
